\documentclass[11pt,a4paper,twoside]{report} % 'twoside' ist optimal für den Druck

% --- Grundlegende Pakete ---
\usepackage[utf8]{inputenc}
\usepackage[ngerman]{babel}
\usepackage[T1]{fontenc}
\usepackage{amsmath, amssymb, amsfonts}
\usepackage{graphicx}
\usepackage{geometry}
\geometry{left=3cm, right=2.5cm, top=3cm, bottom=3.5cm, headheight=15pt}
\usepackage{hyperref}

% --- Kopf- und Fußzeilen (fancyhdr) ---
\usepackage{fancyhdr}
\pagestyle{fancy}
\fancyhf{} % Bereinigt die Standard-Einstellungen
\fancyhead[L]{\nouppercase{\leftmark}} % Kapitelname oben links
\fancyhead[R]{Modellierung komplexer Systeme} % Kurztitel oben rechts
\fancyfoot[L]{R. Fazekas \& A. Makaratzis} % Eure Namen unten links
\fancyfoot[R]{Seite \thepage} % Seitenzahl unten rechts
\renewcommand{\headrulewidth}{0.4pt} % Trennlinie Kopfzeile
\renewcommand{\footrulewidth}{0.4pt} % Trennlinie Fußzeile

% --- Literaturverzeichnis Setup (biblatex) ---
\usepackage[backend=biber,style=ieee]{biblatex}
\addbibresource{references.bib} % Verweist auf die .bib Datei

% --- Metadaten für die Titelseite ---
\title{\textbf{Dokumentation: Modellierung komplexer Systeme}\\ \large Systemarchitektur, Dynamik und Validierung}
\author{Roberto Fazekas \and Alexis Makaratzis}
\date{\today}

\begin{document}

% Titelseite generieren
\maketitle

% Inhaltsverzeichnis
\tableofcontents
\newpage

% --- Hauptteil (Kapitel einbinden) ---
% --- Hauptebene: Kapitel ---
\chapter{Systemarchitektur und Dynamik}
\label{ch:architektur}

Dieses Kapitel beschreibt den grundlegenden Aufbau des Modells und stützt sich auf die theoretischen Konzepte der nichtlinearen Dynamik \cite{strogatz2018}.

% --- Zweite Ebene: Sektion ---
\section{Definition der Subsysteme}
\label{sec:subsysteme}

Das Gesamtsystem wird in drei gekoppelte Subsysteme unterteilt. Wie in Abbildung \ref{fig:regelkreis} zu sehen ist, spielt die Rückkopplung eine zentrale Rolle für die Stabilität.

% --- Einbindung eines Bildes mit Beschreibung ---
\begin{figure}[htbp]
    \centering
    % Ersetze 'regelkreis.png' durch deinen tatsächlichen Dateinamen
    % \includegraphics[width=0.8\textwidth]{images/regelkreis.png}
    
    % \caption[Kurztext fürs Verzeichnis]{Langtext für unter das Bild}
    \caption[Blockschaltbild des gekoppelten Regelkreises]{Detaillierte Darstellung der Rückkopplungsschleifen innerhalb des Hauptsystems. Beachten Sie die Verzögerungsglieder in der Feedback-Schleife, welche die Oszillationen im Systemzustand verursachen.}
    \label{fig:regelkreis}
\end{figure}

% --- Dritte Ebene: Untersektion ---
\subsection{Mathematische Formulierung der Kopplung}

Die Interaktion zwischen den Modulen lässt sich durch folgende Matrixgleichung beschreiben, wobei $K$ die Kopplungsmatrix darstellt:
$$ \dot{X} = F(X) + K \cdot G(X) $$

% \input{sections/02_...}

% --- Verzeichnisse am Ende ---
\newpage
\phantomsection
\addcontentsline{toc}{chapter}{Abbildungsverzeichnis}
\listoffigures % Generiert das Bildverzeichnis automatisch

\newpage
\phantomsection
\addcontentsline{toc}{chapter}{Literaturverzeichnis}
\printbibliography % Generiert das Literaturverzeichnis automatisch

\end{document}
